\documentclass[11pt,a4paper]{article}

\usepackage[margin=1.1in]{geometry}
\usepackage[T1]{fontenc}
\usepackage{lmodern}
\usepackage{amsmath,amssymb}
\usepackage{xcolor}
\usepackage{listings}
\usepackage{booktabs}
\usepackage{hyperref}
\usepackage{microtype}
\usepackage{enumitem}
\setlist{itemsep=2pt,topsep=4pt}

\hypersetup{colorlinks=true,linkcolor=blue!50!black,urlcolor=blue!60!black,citecolor=green!50!black}

\definecolor{codebg}{RGB}{248,248,248}
\definecolor{codeframe}{RGB}{200,200,200}
\definecolor{ckw}{RGB}{0,0,180}
\definecolor{ccom}{RGB}{100,120,100}
\definecolor{cstr}{RGB}{160,60,60}

\lstset{
  basicstyle=\ttfamily\small,
  backgroundcolor=\color{codebg},
  frame=single,rulecolor=\color{codeframe},
  xleftmargin=4pt,xrightmargin=4pt,
  aboveskip=6pt,belowskip=6pt,
  keywordstyle=\color{ckw}\bfseries,
  commentstyle=\color{ccom}\itshape,
  stringstyle=\color{cstr},
  showstringspaces=false,
  columns=fullflexible,keepspaces=true,
}

% Small box for byte-count deltas
\newcommand{\bytes}[1]{\textcolor{blue!60!black}{\textbf{#1\,B}}}
\newcommand{\saved}[1]{\textcolor{green!50!black}{\textbf{$-$#1\,B}}}
\newcommand{\spent}[1]{\textcolor{red!70!black}{\textbf{$+$#1\,B}}}

\title{\bfseries Shrinking an ECDSA Verifier to 890 Bytes\\[0.3em]
       \large A beginner's guide to the tricks that worked}
\author{}
\date{March 2026}

\begin{document}
\maketitle

\begin{abstract}\noindent
This report walks through how an ECDSA/P-256 signature verifier was
shrunk from \bytes{3076} of portable C down to \bytes{890} of x86-64
machine code --- a factor of 3.5$\times$. The target context is a
boot~ROM: the code must fit in a tiny mask-programmed memory, has no
libc, and verifies exactly one signature before jumping to firmware.
Every byte in a boot~ROM is paid for in silicon.

We focus on \emph{why} each trick works, not just \emph{that} it
works. No prior exposure to elliptic curves is assumed; we introduce
just enough of the mathematics to make sense of the size wins. The
full source, test suite, and commit history live in
\texttt{tv\_ecdsa\_tiny.S} and the surrounding repository.
\end{abstract}

\tableofcontents

% ======================================================================
\section{The problem}

A boot~ROM is a small block of read-only memory burned into a chip at
the factory. When the chip powers on, it runs this code first --- before
any firmware, before any operating system. If the boot~ROM verifies
the firmware's signature before handing over control, the chip has a
\emph{hardware root of trust}: even an attacker who can overwrite
flash cannot run unsigned code.

The catch is size. Mask~ROM cells are large compared to flash, every
chip on the wafer carries the same bits, and respinning silicon to fix
a bug costs months and seven figures. Manufacturers budget boot~ROMs in
\emph{kilobytes}. Your verifier shares that budget with the flash
driver, the clock setup, and whatever else the chip needs to come
alive. If you spend 3\,KB on signature verification, something else
gets cut.

So the game is: \textbf{implement ECDSA signature verification over
curve P-256, as small as you can make it, without cutting corners on
correctness.}

\subsection{The ground rules}

\begin{itemize}
\item \textbf{Count everything.} Size is measured on the \emph{object
  file} --- text plus read-only data. If you call \texttt{memcpy},
  \texttt{memcpy} counts. The target here has zero undefined symbols:
  no libc, no compiler runtime, nothing.
\item \textbf{Correctness is not negotiable.} Every build must pass
  607 test vectors: 33 hand-picked edge cases (zero components,
  out-of-range values, the point at infinity, signature malleability)
  plus 574 from Google's Wycheproof suite~\cite{wycheproof}. A
  verifier that's ten bytes smaller but accepts a forged signature is
  worth exactly nothing.
\item \textbf{Speed is secondary but not irrelevant.} A boot~ROM
  runs once per power-on, so milliseconds are fine. But there's a
  \emph{Pareto frontier}: if build~A is both smaller and
  faster than build~B, build~B is dead. We track both axes.
\item \textbf{Constant time is not required.} Signature
  \emph{verification} takes only public inputs --- the signature, the
  public key, the message hash. Nothing secret. Data-dependent
  branches and variable-time reductions, which would be disqualifying
  in a \emph{signer}, are fair game here. Several of the tricks below
  exploit this freedom.
\end{itemize}

\subsection{What ECDSA verification actually computes}
\label{sec:ecdsa}

You do not need to know why ECDSA is secure to follow this report.
You only need to know what the verifier does, because that
determines where the bytes go.

An ECDSA signature over P-256 is a pair of integers $(r, s)$, each
256 bits wide. The public key is a point $Q$ on an elliptic curve.
Verification, per FIPS~186-5~\cite{fips186}, computes:
\begin{align*}
  w   &= s^{-1} \bmod n            &&\text{\small (one modular inversion)}\\
  u_1 &= e \cdot w \bmod n         &&\text{\small ($e$ = the message hash)}\\
  u_2 &= r \cdot w \bmod n         \\
  (x, y) &= u_1 \cdot G + u_2 \cdot Q  &&\text{\small (two scalar multiplications, one point add)}\\
  &\text{accept iff } x \equiv r \pmod{n}
\end{align*}
Here $G$ is a fixed ``generator'' point baked into the standard, $n$
is the order of~$G$ (a 256-bit prime), and $p$ is the field modulus
(another 256-bit prime). All the elliptic-curve arithmetic happens
modulo~$p$; the scalars $u_1, u_2, w$ live modulo~$n$.

The heavy lifting is in $u_1 \cdot G + u_2 \cdot Q$. Computing
$k \cdot P$ means ``add the point $P$ to itself $k$ times'' --- but
since $k$ is a 256-bit number, you do this with a square-and-multiply
style loop: 256 iterations of \emph{point doubling} (compute $2P$
from $P$) interleaved with \emph{point addition} (compute $P+Q$ from
$P$ and $Q$) wherever the scalar has a 1 bit.

Each point doubling or point addition is in turn a fixed sequence of
field operations: multiplications, squarings, additions, subtractions,
all modulo~$p$. A P-256 point addition is roughly a dozen field
multiplications plus change. One standard trick: instead of working
with $(x, y)$ pairs directly, the loop carries an extra coordinate
$Z$ and represents the running point as $(X, Y, Z)$ with $x = X/Z$.
Divisions are expensive (see \S\ref{sec:projcheck}); this
\emph{projective} representation defers them all to one division at
the very end.

\textbf{So the byte budget is roughly:} one field multiplier (mod~$p$
and mod~$n$), one point-add and one point-double formula that call it,
one scalar-multiply loop that calls them, one inversion, and
the top-level glue that decodes the inputs and checks the result.
Constants for $p$, $n$, $G$, and the curve coefficient $b$ take
another 128--160 bytes.

% ======================================================================
\section{The journey at a glance}

Table~\ref{tab:journey} gives the headline numbers. The rest of this
document goes through each row in enough detail that you could, in
principle, rediscover the trick yourself.

\begin{table}[h]
\centering
\begin{tabular}{@{}rrl@{}}
\toprule
Size & $\Delta$ & What changed \\
\midrule
3076\,B &          & Portable C, \texttt{gcc -Os}, 32-bit limbs \\
2873\,B &          & First hand-written x86-64 assembly \\
1712\,B & $-$1161  & \S\ref{sec:bytecode}: \textbf{Bytecode interpreter} \\
1427\,B & $-$285   & \texttt{fast.S}: Montgomery + \texttt{mulx} + Shamir's trick \\
1397\,B & $-$30    & Micro-grind (register allocation, fall-through) \\
\midrule
1300\,B & $-$97    & \S\ref{sec:nomont}: \textbf{Drop Montgomery --- the $q=t[\text{top}]$ reduce} \\
1195\,B & $-$105   & \S\ref{sec:projcheck}: \textbf{Projective final check --- no mod-$p$ inversion} \\
1105\,B & $-$90    & \S\ref{sec:btcn}: \textbf{\texttt{bt} on $n$ directly --- no exponent buffer} \\
1012\,B & $-$93    & \S\ref{sec:rcb}: \textbf{RCB complete addition --- one formula, no branches} \\
 985\,B & $-$27    & \S\ref{sec:slotshift}: \textbf{Addend slot shift --- layout is free, setup is not} \\
 942\,B & $-$43    & \S\ref{sec:runtimecp}: \textbf{Build $p$ at runtime} + fe\_iszero inlines \\
 933\,B & $-$9     & \S\ref{sec:efd}: \textbf{EFD reschedule} --- reorder for disp8 reach \\
 898\,B & $-$35    & \S\ref{sec:below}: \textbf{Over-copy, decoder preload, paired decode} \\
 890\,B & $-$8     & \S\ref{sec:grind}: \texttt{repe scasq} zero-check, suffix-sharing \\
\bottomrule
\end{tabular}
\caption{Major milestones. Sizes between rows are connected by dozens
  of smaller wins; see \S\ref{sec:grind} for the general flavour. The
  rule marks where \texttt{tiny.S} forked from \texttt{fast.S}: rows
  above it optimised for speed first; rows below trade speed for
  size.}
\label{tab:journey}
\end{table}

% ======================================================================
\section{Bytecode: call sites are expensive, interpreters are cheap}
\label{sec:bytecode}

\noindent\emph{From \bytes{2873} to \bytes{1712}.}

\subsection{The problem with calling functions}

The first hand-written assembly version looked the way you'd expect:
a function for field-multiply, a function for field-add, a function
for field-subtract, and then \texttt{pt\_dbl} and \texttt{pt\_add}
calling them in sequence.

On x86-64, calling a three-argument function costs about 17 bytes:
\begin{lstlisting}[language={[x86masm]Assembler}]
    lea   rdi, [r14 + 32*DST]    ; 4 bytes (or 7 if disp > 127)
    lea   rsi, [r14 + 32*SRC1]   ; 4 bytes
    lea   rdx, [r14 + 32*SRC2]   ; 4 bytes
    call  fe_mul                 ; 5 bytes
\end{lstlisting}
Point doubling is about 8 field operations; point addition is about
12. The verifier's final-check sequence is another 10 or so. That's
roughly 60 call sites~$\times$~15~bytes each: \textbf{900~bytes spent
on nothing but ``load three pointers and jump.''}

\subsection{The fix: a tiny virtual machine}

Notice that every one of those call sites has the same \emph{shape}:
one destination, two sources, one operation. And all the operands
live in a contiguous block of 32-byte field elements on the stack.
So instead of emitting 15 bytes of setup per call, we encode each
call as \textbf{2 bytes of data} and write a small loop that decodes
and dispatches:

\begin{lstlisting}[language={[x86masm]Assembler}]
    ; Each instruction: byte0 = (s2<<4)|op, byte1 = (dst<<4)|s1
    ; With 16 slots and 16 ops, everything fits in nibbles.
bc_run:
    lodsw                        ; fetch 2 bytes of bytecode into ax
    ; ... decode nibbles into rdi/rsi/rdx (about 30 bytes) ...
    movzx ecx, byte [r12 + rax]  ; jump-table lookup (op index -> handler offset)
    add   rcx, r12
    call  rcx                    ; dispatch
    jmp   bc_run
\end{lstlisting}

The dispatch loop is about 60 bytes. Each handler (multiply, add,
subtract, \ldots) is another 10--40 bytes --- but you already had
those as functions. The bytecode streams themselves are $60 \times 2
= 120$ bytes.

\textbf{Net effect:} $\sim$900 bytes of call sites becomes $\sim$180
bytes of interpreter plus bytecode. The \texttt{pt\_dbl} function
literally \emph{disappears} as native code --- it is now 25 bytes of
data in \texttt{.rodata}.

\subsection{Why 2 bytes, why nibbles}

You could encode each op in 3 bytes (one per operand) and have room
for 256 slots. But you don't have 256 slots --- point addition
uses perhaps a dozen temporaries. Four bits per operand, packed two to
a byte, is the right granularity. And four bits of opcode gives you 16
operations --- the final build uses ten.

There's a subtle win hiding here: \textbf{the interpreter doesn't
care which \emph{specific} 16 slots you use}. You can reassign slots
to make other code shorter. We will exploit this ruthlessly later
(\S\ref{sec:slotshift}).

% ======================================================================
\section{Dropping Montgomery: the $q = t[\text{top}]$ reduce}
\label{sec:nomont}

\noindent\emph{From \bytes{1397} to \bytes{1300}. Cornerstone of the
size-optimised build.}

\subsection{What Montgomery multiplication costs you}

The fast implementation used \emph{Montgomery form}~\cite{montgomery}:
numbers are stored as $a \cdot R \bmod m$ where $R = 2^{256}$. This
makes the reduce-after-multiply step very fast --- but it carries
hidden fixed costs:
\begin{itemize}
\item A precomputed constant $m_0^{-1} = -m^{-1} \bmod 2^{64}$ for
  each modulus ($p$ and $n$): \bytes{16} of data plus the code to
  pass it into the multiplier on every call.
\item A precomputed constant $R^2 \bmod p$ to convert inputs
  \emph{into} Montgomery form: \bytes{32} of data plus a
  \texttt{TOMONT} bytecode op.
\item Conversion \emph{out} of Montgomery form at the end: more
  bytecode ops.
\end{itemize}
None of these are big individually, but they add up to about 60--80
bytes --- and for size, we'd rather have a slower multiplier with no
baggage.

\subsection{The special shape of the P-256 moduli}

Here is where we get lucky. Write out $p$ and $n$ in hex and stare at
the top word:
\begin{align*}
p &= \texttt{ffffffff\,00000001\,00000000\,00000000\,00000000\,ffffffff\,ffffffff\,ffffffff} \\
n &= \texttt{ffffffff\,00000000\,ffffffff\,ffffffff\,bce6faad\,a7179e84\,f3b9cac2\,fc632551}
\end{align*}
\textbf{Both moduli have a top 32-bit word of \texttt{0xFFFFFFFF}.}
And both satisfy $2^{256} - m < 2^{224}$ --- that is, both are ``just
barely less than'' $2^{256}$.

\subsection{Why that shape makes reduction trivial}

After a schoolbook multiply, you have a 512-bit product $t$ and you
want $t \bmod m$. The textbook approach is trial division: estimate
the quotient, multiply, subtract, correct. Quotient estimation is
where the complexity lives.

But if the top word of $m$ is \texttt{0xFFFFFFFF}, quotient estimation
becomes \textbf{free}. Split $t$ into 32-bit words $t[0..15]$ and
slide a 9-word window from the top down. At each position $j$:
\begin{enumerate}
\item Let $q = t[j{+}8]$ --- just \emph{read} the top word of the
  window. No division.
\item Subtract $q \cdot m$ from the window.
\end{enumerate}

Why is $q = t[j{+}8]$ a good quotient estimate? Because $m$'s top
word is $2^{32}-1$, the true quotient
$\lfloor t_{\text{window}} / m \rfloor$ (the brackets mean round down) is within 1 of
$\lfloor t_{\text{window}} / 2^{256} \rfloor = t[j{+}8]$.

Concretely: write the window as $W = q \cdot 2^{256} + L$ with
$0 \le L < 2^{256}$. Then
\[
  W - q \cdot m \;=\; q\,(2^{256} - m) + L
  \;<\; 2^{32} \cdot 2^{224} + 2^{256}
  \;=\; 2 \cdot 2^{256},
\]
so after one subtraction the window is at most two moduli wide, and
its new top word --- $\lfloor (W - qm) / 2^{256} \rfloor$ --- is
either 0 or 1. If it's~1, one more pass brings it to~0. Two passes
per position, worst case.

\begin{lstlisting}[language={[x86masm]Assembler}]
; Reduce: slide a 9-dword window from j=8 down to j=0.
; At each j: q = t[j+8];  t[j..j+8] -= q*m.
; The SAME 26-byte primitive handles both schoolbook multiply
; AND this reduce --- they're both "mul ebx; accumulate into [rdi]".
.Lreduce:
    mov   ebx, [rdi + 32]   ; q = top dword of current window
    ; ... call the shared mul-accumulate primitive ...
    sub   rdi, 4            ; slide window down one dword
    dec   ecx
    jnz   .Lreduce
\end{lstlisting}

\textbf{The payoff:} no $m_0^{-1}$, no $R^2$, no conversion ops.
Field elements are stored in plain form the whole time. The reduce
code is \emph{shorter} than Montgomery's, and one 26-byte inner
primitive serves double duty for multiply and reduce.

The trade-off is speed: this runs at 32-bit granularity ($8 \times 8$
word multiplies instead of $4 \times 4$), so it's about $2\times$
slower. For a boot~ROM, we'll take that trade every time.

\paragraph{A subtlety worth noting.} The $q = t[\text{top}]$ trick
requires the top \emph{32-bit} word to be all-ones. The top
\emph{64-bit} word of $p$ is \texttt{0xFFFFFFFF00000001}, which is
\emph{not} all-ones. So even after we later upgraded the schoolbook
multiply to 64-bit limbs (\spent{14} for about a $2\times$ speedup),
the reduce stayed at 32-bit granularity. This is why
\texttt{tv\_ecdsa\_tiny.S} has a slightly odd split: 64-bit multiply,
32-bit reduce.

% ======================================================================
\section{The projective final check: deleting an inversion}
\label{sec:projcheck}

\noindent\emph{From \bytes{1216} to \bytes{1195}. (Table~\ref{tab:journey}
attributes \bytes{105} to this milestone: the gap from \bytes{1300} is
micro-grind between milestones; this section covers the final idea.)
One of those rare wins that's both smaller and faster.}

\subsection{Where the mod-$p$ inversion was hiding}

The ECDSA check (last line of \S\ref{sec:ecdsa}) is ``accept if
$x \equiv r \pmod n$,'' where $x$ is the affine $x$-coordinate of the
computed point. But the scalar multiplication loop doesn't produce
affine coordinates --- it produces \emph{projective} coordinates
$(X, Y, Z)$, where the affine $x$ is recovered as $x = X / Z$ (for
homogeneous projective) or $x = X / Z^2$ (for Jacobian).

``Divide by $Z$'' in a finite field means ``multiply by $Z^{-1}$.''
For prime~$p$, Fermat's little theorem gives $Z^{p-1} \equiv 1
\pmod p$, so $Z^{-1} \equiv Z^{p-2}$: inversion is exponentiation.
That exponentiation is a 256-iteration square-and-multiply loop ---
another $\sim$380 field multiplications on top of everything else.

The straightforward verifier does this inversion, recovers $x$, and
compares it to $r$. That inversion call, its argument setup, and the
final compare-and-reduce-mod-$n$ block cost about 45 bytes.

\subsection{Clearing denominators}

Here's the algebra trick. Instead of computing $x = X/Z$ and checking
$x \stackrel{?}{=} r$, \textbf{clear the denominator}: check
$X \stackrel{?}{\equiv} r \cdot Z \pmod p$. No inversion needed.

There's one wrinkle. The ECDSA check is modulo $n$, but $X$ and $Z$
live modulo $p$, and $p \neq n$. Since $0 \le x < p < 2n$, there are
at most two integers congruent to $r \pmod n$ in that range: $r$
itself, and $r + n$ (the latter only when $r + n < p$). So the full
check becomes:
\[
  X \equiv r \cdot Z \pmod p
  \quad\textsc{or}\quad
  X \equiv (r+n) \cdot Z \pmod p.
\]
When $r < p - n$ --- a roughly $2^{-128}$ sliver of the $r$ range,
since $p - n$ is only a 127-bit number --- the second disjunct is
genuinely needed. The rest of the time $r + n$ wraps around mod~$p$
and the second disjunct tests a spurious value. That's harmless:
forcing the scalar multiplication to land on a \emph{specific}
spurious point is as hard as the discrete log problem (the hard
mathematical problem ECDSA's security rests on).

\subsection{Folding the \textsc{or} into a multiply}

Checking an \textsc{or} in assembly usually means a branch, and
branches are byte-expensive. But we're working modulo a prime $p$,
and in a field, \textbf{a product is zero iff one of its factors is
zero}. So let
\[
  d_1 = X - r\cdot Z, \qquad d_2 = X - (r+n)\cdot Z,
\]
and check whether $d_1 \cdot d_2 \equiv 0 \pmod p$. One multiply,
one zero-check, no branches.

The whole thing --- compute $r\cdot Z$, compute $d_1$, compute $d_2$,
multiply, check-zero --- is 7 field operations, which is 14 bytes of
bytecode. The 45 bytes of native inversion-and-compare are gone,
replaced by 14 bytes of data.

And as a bonus, we no longer call the inverter with modulus $p$ at
all. The inverter now only ever runs mod~$n$ (for $s^{-1}$). That
removes an \texttt{INV\_P} handler from the bytecode jump table and
lets us hardcode $n$ into the inverter, saving another handful of
bytes (\S\ref{sec:btcn}).

% ======================================================================
\section{\texttt{bt} on $n$ directly: no exponent buffer}
\label{sec:btcn}

\noindent\emph{From \bytes{1124} to \bytes{1105}.}

\subsection{The cost of computing $n - 2$}

Fermat's little theorem gives $s^{-1} \equiv s^{n-2} \pmod n$. The
exponentiation is a 256-step square-and-multiply loop that reads one
bit of the exponent $n-2$ per iteration: square always, multiply if
the bit is 1.

The exponent is $n - 2$. Previously we stored $n$ as a constant, and
at runtime: copy $n$ to a 32-byte stack buffer, subtract 2 from the
buffer, then have the loop read bits from the buffer with the
\texttt{bt} instruction.

That's an \texttt{enter} to allocate the buffer, a \texttt{rep movsq}
to copy, a \texttt{sub} to decrement, and a \texttt{leave} to clean
up: 33 bytes of machinery just to turn $n$ into $n-2$.

\subsection{The bit-level observation}

Write out the low byte of $n$ and $n-2$ in binary:
\begin{center}
\begin{tabular}{r@{\;=\;}l@{\;=\;}l}
$n \bmod 256$   & \texttt{0x51} & \texttt{0101\,0001} \\
$(n-2) \bmod 256$ & \texttt{0x4F} & \texttt{0100\,1111} \\
XOR               & \texttt{0x1E} & \texttt{0001\,1110} \\
\end{tabular}
\end{center}
Subtracting 2 borrows from bit~1. That bit is 0 in~$n$, so the borrow
propagates: through bit~2 (also~0), through bit~3 (also~0), until it
reaches bit~4 --- which is~1 and absorbs it. \textbf{Bits 5 through
255 of $n-2$ are identical to bits 5 through 255 of $n$.}

So: don't compute $n-2$ at all. Read bits directly from the $n$
constant that's already sitting in \texttt{.rodata}. For bits 5 and
up, $n$ gives the right answer. For bits 0--4, patch in a 7-byte
special case:
\begin{lstlisting}[language={[x86masm]Assembler}]
    cmp   ebx, 4
    jb    .Lmul       ; bits 0-3: n-2 has 1111 here, force multiply
    je    .Lnext      ; bit 4:    n-2 has 0, n has 1, force skip
    bt    [r12 + oN], ebx   ; bits 5+: n and n-2 agree, read n directly
    jnc   .Lnext
.Lmul:
    ; ... multiply step ...
\end{lstlisting}

The 33-byte buffer machinery becomes 7 bytes of branching. Six bytes
of the original machinery survive (the init copy still has to run),
and the \texttt{bt} operand grows by one byte (it now reaches into
\texttt{.rodata} instead of \texttt{[rsp]}). Net: $-33 + 7 + 6 + 1 =
\saved{19}$.

\subsection{The broader lesson}

This trick is \emph{ludicrously specific} to the bit pattern of the
P-256 group order. It would not work for most other primes. But
that's exactly the point: when optimising for size, there is no
``generic'' --- you stare at the \emph{actual} constants and ask what
accidents of their structure you can exploit.

% ======================================================================
\section{RCB complete addition: one formula for everything}
\label{sec:rcb}

\noindent\emph{From \bytes{1071} to \bytes{1012}. The biggest
single-commit win after the bytecode interpreter.}

\subsection{Why point addition normally has three cases}

The classical formulas for adding two elliptic curve points $P + Q$
have a dirty secret: they divide by $x_P - x_Q$. When $P = Q$
(doubling), that's $0/0$. When $P = -Q$ (the result is the point at
infinity $\mathcal{O}$), same problem. And if either input is already
$\mathcal{O}$, the formulas don't even apply --- $\mathcal{O}$ has no
coordinates.

So a traditional \texttt{pt\_add\_acc} routine starts with a
three-way branch:
\begin{lstlisting}[language=C]
  if (acc.z == 0)     { acc = Q; return; }     // acc was infinity
  // ... run first half of add formula to compute H = x_Q - x_P ...
  if (H == 0) {
      if (Rr == 0)    { pt_dbl(acc); return; } // P == Q, use doubling
      else            { acc = infinity; return; } // P == -Q
  }
  // ... finish the addition formula ...
\end{lstlisting}
In our bytecode world this was 62 bytes of native branching, plus
\emph{three} separate bytecode streams: \texttt{bc\_dbl} (50\,B),
\texttt{bc\_add1} (24\,B, the pre-check half), and \texttt{bc\_add2}
(28\,B, the post-check half). Total: 164 bytes.

\subsection{The complete formula}

In 2015, Renes, Costello, and Batina~\cite{rcb} published a point
addition formula that \textbf{has no exceptional cases}. One sequence
of 43 field operations gives the correct answer for \emph{every} pair
of inputs: distinct points, equal points, inverse points, and the
point at infinity. No branches. No special cases. It just works.

The formula uses \emph{homogeneous} projective coordinates: the
affine point is recovered as $(X/Z,\, Y/Z)$. Two triples name the
same point
whenever they differ by a common nonzero scale factor: $(X:Y:Z)$ and
$(\lambda X:\lambda Y:\lambda Z)$ are the same thing, which is what
the colons signal. Infinity is
the class with $Z = 0$, which on this curve forces $X = 0$ too: any
$(0:Y:0)$ with $Y \neq 0$ will do. Feeding the formula $(0:1:0) + Q$
really does produce~$Q$, and $(0:1:0) + (0:1:0)$ really does produce
$(0:1:0)$ --- the edge cases are handled by the algebra, not by code.

\subsection{The demolition}

With RCB, \texttt{pt\_add\_acc} becomes:
\begin{lstlisting}[language={[x86masm]Assembler}]
pt_add_acc:                 ; 11 bytes.  Was 62.
    ; Addend already in slots 5,6,7.  Just dispatch bc_rcb.
    xor   edi, edi          ; bytecode offset 0 = bc_rcb
    jmp   bcrun_off
\end{lstlisting}
Doubling is not a separate operation any more. To double the
accumulator, copy it into the addend slots and call the \emph{same}
routine --- $P + P = 2P$, and RCB handles that case like any other.

The three bytecode streams (\bytes{102} total) collapse into one
43-operation stream (\bytes{88}). The 62-byte branch tree becomes
an 11-byte dispatch. Setup bookkeeping for the new $Z$ coordinate
eats \spent{6} of that back; net \saved{59}, and a whole category
of bugs (forget to handle one edge case, accept a forged
signature) structurally eliminated.

\subsection{The fine print}

RCB does more arithmetic than the classical ``generic add'' --- 43
operations versus about 16. So each addition is about $2.5\times$
slower. For a boot~ROM, we shrug. For a speed-optimised build, you'd
keep the branchy version.

One invariant turned out to be load-bearing: the RCB formula never
\emph{writes} to the slots holding $Q$. It reads them, but the
outputs go elsewhere. This means the scalar-multiply loop doesn't
need to save and restore $Q$ across each add --- $Q$ just sits there,
undisturbed, for all 256 iterations. That invariant saved another
dozen bytes of copy-in-copy-out in the loop.

% ======================================================================
\section{Layout is free, setup is not}
\label{sec:slotshift}

\noindent\emph{From \bytes{1001} to \bytes{985}.}

The bytecode interpreter addresses 16 slots by a 4-bit index. Which
slot is ``slot~5'' is completely arbitrary --- it's a nibble in a
data table. Renumbering slots costs \textbf{zero bytes} in the
bytecode. But slot numbers show up in \emph{native} code too, as
memory displacements --- and those are not free.

\subsection{Shamir's trick and its setup cost}

Shamir's trick~\cite[\S 14.88]{hac} (due originally to
Straus~\cite{straus}) computes $u_1 G + u_2 Q$ with a single
256-iteration loop instead of two: at each bit
position, double once, then conditionally add $G$, conditionally
add~$Q$. It halves the number of doublings.

The trick needs backup copies of both $G$ and $Q$ that survive the
whole loop, because the ``working'' addend slots get overwritten.
After the RCB switch, each point is 3 slots $(X, Y, Z)$. So the
setup is: copy $G$ (3 slots) and $Q$ (3 slots) to backup slots
16--21. (These live outside the bytecode's 4-bit addressable range
--- only native code touches them.)

The naive setup was three separate copies: $G$'s $X,Y$ from one
place, a borrowed $Z=1$ from another, $Q$ from a third. About 28
bytes of \texttt{lea}/\texttt{rep movsq} triples.

\subsection{Making the source contiguous}

What if the six source slots were \emph{already contiguous} before
the copy? Then the setup is one \texttt{rep movsq} with count 24.
Ten bytes instead of 28.

To make this happen: (a) arrange for the early validation bytecode
to leave slots 2--7 holding exactly $G_x, G_y, 1, Q_x, Q_y, 1$ in
that order; (b) shift the RCB addend from slots 4,5,6 to slots
5,6,7 so it lines up.

Part (a) costs 2 extra bytes of bytecode (one more ``copy 1 here''
op). Part (b) is a \textbf{pure nibble permutation}: swap some
4-bit fields in \texttt{bc\_rcb}, zero-byte cost. Net: \saved{16}.

The general principle: data layout changes inside the bytecode are
free; exploit that freedom to make the native code shorter.

% ======================================================================
\section{Building constants at runtime}
\label{sec:runtimecp}

\noindent\emph{From \bytes{948} to \bytes{942} --- but the direct
saving is zero. The entire win is a knock-on effect.}

\subsection{A constant you can compute in the same number of bytes}

The field prime $p$ for P-256 is
\[
  p = 2^{256} - 2^{224} + 2^{192} + 2^{96} - 1.
\]
As 32 bytes of \texttt{.rodata} it's \bytes{32}. But look at the hex:
\[
\texttt{ffffffff\,00000001\,00000000\,00000000\,00000000\,ffffffff\,ffffffff\,ffffffff}
\]
Seven of the eight 32-bit words are either \texttt{0x00000000} or
\texttt{0xFFFFFFFF} --- only \texttt{0x00000001} is neither. Writing
this out one dword at a time with \texttt{stosd} (which stores
\texttt{eax} and auto-advances \texttt{rdi}) takes \bytes{17}:
\begin{lstlisting}[language={[x86masm]Assembler}]
    or    eax, -1       ; eax = FFFFFFFF   (3 B)
    stosd               ; dword 0          (1 B each; stosd = AB)
    stosd               ; dword 1
    stosd               ; dword 2
    xor   eax, eax      ; eax = 0          (2 B)
    stosd               ; dword 3
    stosd               ; dword 4
    stosd               ; dword 5
    inc   eax           ; eax = 1          (2 B)
    stosd               ; dword 6
    neg   eax           ; eax = FFFFFFFF   (2 B)
    stosd               ; dword 7
\end{lstlisting}
That looks like a \saved{15} win --- but the constant didn't sit in a
vacuum. Deleting it means the places that \emph{used} it need new
plumbing: a pointer register to hold the built-at-runtime address, a
new bytecode op so validation can check coordinates against it,
save/restore of the new register in the prologue. About \spent{15} of
overhead, netting out to roughly zero.

So why bother? Because \textbf{where the 32 bytes disappeared from
matters more than how many were saved.}

\subsection{The knock-on win}

The bytecode jump table stores handler addresses as \emph{8-bit}
offsets (one byte per handler). This means every handler must be
within 255 bytes of the table base. The largest handler
(\texttt{fe\_inv\_m}) was sitting at offset 247 --- only 8 bytes of
headroom before the whole scheme breaks and every table entry doubles
in size.

The \texttt{cP} constant lived between the jump table and the
handlers. Deleting it from \texttt{.rodata} moved every handler 32
bytes closer to the table. Suddenly there's 36 bytes of headroom ---
enough to inline a small \texttt{fe\_iszero} helper into its only
caller, for \saved{6}. That is the entire saving in the section
header.

(We had already done the related trick for the curve coefficient $b$,
computed \emph{in bytecode} from the curve equation itself: since $G$
is on the curve, $b = G_y^2 - G_x^3 + 3 G_x \pmod p$. Four field ops,
8 bytes of bytecode, replacing \bytes{32} of constant --- a direct
saving this time, because bytecode costs no plumbing.)

Size optimisations interact: freeing space in one region can
unblock a seemingly unrelated win elsewhere. Keep a list
of ``things I'd do if only I had 10 more bytes of headroom at $X$.''

% ======================================================================
\section{The EFD reschedule: reorder for displacement reach}
\label{sec:efd}

\noindent\emph{From \bytes{935} to \bytes{933}.}

x86-64 memory operands with a small displacement ($-128 \le d \le
127$) encode in one byte; larger displacements take four. One
\texttt{lea rdi, [r14+160]} in the hot copy path was paying the
three-byte tax because slot~5 is $5 \times 32 = 160$ bytes from the
slot base, and $160 > 127$.

If we had a register that already pointed near slot~5, we could
use a short displacement from that register. And \texttt{r8} was
already holding \texttt{\&cP} --- but \texttt{cP} was at slot~22, way
too far.

\textbf{So: move \texttt{cP}.} If \texttt{cP} lived at slot~8 instead
($8 \times 32 = 256$ from base), then slot~5 would be at
\texttt{r8 $-$ 96}, which fits in a signed byte. But slot~8 was being
used as scratch by the RCB formula.

The Explicit-Formulas Database~\cite{efd} lists the RCB operations in
a particular order, but the order is not sacred --- only the
\emph{data dependencies} are. By hoisting two independent steps
earlier in the sequence, $Z_1$ dies sooner, slot~2 becomes reusable
as a sixth temporary, and slot~8 is \emph{never written} by
\texttt{bc\_rcb}. It can hold \texttt{cP} for the entire scalar
multiplication.

Shuffling bytecode nibbles is free. The displacement shrinks from
\texttt{+160} (4 bytes) to \texttt{$-$96} (1 byte). Minus one byte
for a setup adjustment elsewhere, net \saved{2}.

Two bytes is small, but it illustrates how far the rabbit hole goes:
\emph{instruction scheduling in a data table to change register
allocation to shrink an addressing mode in unrelated native code.}

% ======================================================================
\section{Below what looked like the floor}
\label{sec:below}

\noindent\emph{From \bytes{933} to \bytes{890}.}

The 933-byte build sat unchanged for long enough that it started to
feel like a hard limit for this architecture. It wasn't. None of the
following changed the algorithm --- still the same $q=t[\text{top}]$
reduce, same RCB formula, same bytecode --- but 35 more bytes came
out of the encoding and the layout.

\subsection{Over-copy: let \texttt{rep movsq} run long}

Several places copy a 3-slot point (12 quadwords) with \texttt{rep
movsq} and then immediately need \texttt{rdi} pointing somewhere
further along for the next operation. The old pattern: copy 12
quadwords, then \texttt{lea rdi, [\ldots]} to reposition --- 4 to 7
bytes of address arithmetic.

Since \texttt{rep movsq} advances both \texttt{rsi} and \texttt{rdi}
by the count, copying \emph{more} than needed does the repositioning
for free. If the slots between the real destination and the next
needed position are scratch --- written before they're next read ---
they can soak up the extra writes without consequence. Doubling a
12-qword copy to 24 qwords costs nothing (\texttt{mov cl,24} is the
same 2 bytes as \texttt{mov cl,12}) and deletes the \texttt{lea}.

This is the dataflow equivalent of writing past the end of your
buffer on purpose, because you own the memory on the other side and
know nobody's looking.

\subsection{Let the decoder do the handler's work}

The bytecode dispatch loop (\S\ref{sec:bytecode}) already does
arithmetic on every fetch: it unpacks three nibbles and turns each
into a slot pointer via \texttt{base + nibble$\times$32}. That
multiply-and-add runs regardless of which opcode follows.

Some handlers were doing their own address arithmetic on entry ---
computing a secondary pointer from a primary one, or offsetting into
a structure. But a slot nibble is four free bits, and a slot pointer
is a free multiply-by-32. If you can encode the handler's desired
offset as a slot number, the decoder has already computed it by the
time the handler runs. The handler shrinks; the bytecode byte it
reads from was already there.

The pattern is general: whenever a fixed-cost decode stage sits in
front of variable handlers, ask what the decode stage could compute
on the handlers' behalf.

\subsection{Paired big-endian decode}

The verifier has to parse four 32-byte big-endian integers from the
input blob: $r$ and $s$ from the signature, $Q_x$ and $Q_y$ from the
public key. They arrive in pairs --- $(r,s)$ adjacent in one buffer,
$(Q_x,Q_y)$ adjacent in another --- and land in adjacent slots.

Each decode was a 5-byte \texttt{call fe\_from\_be}, so each pair
cost 10 bytes of call sites. A 5-byte helper that calls
\texttt{fe\_from\_be} once and then falls through into it a second
time turns each pair into one 5-byte call. Two pairs: \spent{5} for
the helper, \saved{10} at the call sites, net \saved{5}.

% ======================================================================
\section{The micro-grind: x86 encoding folklore}
\label{sec:grind}

Between the big structural wins, hundreds of bytes came out one or
two at a time. These don't make for a good story individually, but
the patterns repeat.

\subsection{\texttt{push;\,pop} is a 2-byte \texttt{mov}}

\texttt{mov rax, rbx} is 3 bytes. \texttt{push rbx; pop rax} is 2
bytes (for registers that don't need a REX prefix). It clobbers 8
bytes of stack below \texttt{rsp}, which is fine if you're not using
the red zone for anything else.

This pattern appears about a dozen times in the final code. It saved
roughly \saved{10} total.

\subsection{Short jumps versus near jumps}

A jump to a target within $\pm 127$ bytes encodes as 2 bytes
(\texttt{EB xx}). Beyond that, it's 5 bytes (\texttt{E9 xx xx xx
xx}). Every function-layout decision --- what order the handlers go
in, whether a helper sits before or after its caller --- affects
which jumps are short.

Several commits in the history are just ``swap these two blocks so
this \texttt{jmp} goes from 5 bytes to 2.'' Three bytes per swap,
maybe five such swaps over the project: \saved{15}.

\subsection{Fall-through instead of jump}

Even better than a short jump is no jump. If function~A always
ends by calling function~B, and nothing else sits between them,
delete the \texttt{call} from~A and let execution fall through.
\texttt{Fadd} falls through into the conditional-subtract epilogue;
\texttt{Nmul} overwrites the modulus pointer and falls through into
\texttt{Fmul}. Each fall-through is \saved{2} to \saved{5}; four or
five over the project comes to roughly \saved{15}.

\subsection{\texttt{rep movsq} is a free \texttt{memcpy}}

\texttt{rep movsq} is 3 bytes and copies \texttt{rcx} quadwords from
\texttt{[rsi]} to \texttt{[rdi]}. If you can get \texttt{rsi},
\texttt{rdi}, and \texttt{rcx} into the right state as a side
effect of surrounding code --- \texttt{lodsb} leaves
\texttt{rsi} advanced, \texttt{stosq} leaves \texttt{rdi} advanced,
loop counters naturally end at zero --- then a block copy is nearly
free.

The final code has about eight \texttt{rep movsq} instructions and
not a single explicit copy loop. The art is in arranging the dataflow
so the registers are already right.

\subsection{Microcode is slow but short}

The \texttt{loop} instruction (decrement \texttt{rcx} and jump if
nonzero) is 2 bytes. The equivalent \texttt{dec rcx; jnz} is 5 bytes.
But \texttt{loop} is \emph{microcoded} on modern x86 and costs
several extra cycles per iteration --- Agner Fog's
tables~\cite{agner} put it at 6 on Skylake, worse on AMD Zen.

If the loop runs ten times, nobody cares. If it's the inner loop of
the field multiplier and runs 950{,}000 times per signature, that's
millions of wasted cycles. The \texttt{-DSMALL\_MUL8} build takes the
size win; the default build spends \spent{14} to run at full speed.
Both are on the Pareto frontier.

\subsection{The \texttt{rcx = 0} flow}

\texttt{mov cl, 8} is 2 bytes and sets the low 8 bits of \texttt{rcx}.
If the upper 56 bits happen to already be zero --- because a previous
\texttt{rep} instruction counted down to zero, or a previous
\texttt{loop} exited --- you just saved 3 bytes over
\texttt{mov ecx, 8} (5 bytes) or 1 byte over \texttt{push 8; pop rcx}
(3 bytes).

This is a \emph{global} invariant: ``\texttt{rcx} is zero at these
twelve points in the control flow, and the code downstream relies on
it.'' It's fragile --- reordering two blocks can silently break it
--- and it saved perhaps \saved{8} across the whole binary. The
source has comments at every point where the invariant is consumed.

% ======================================================================
\section{Things that didn't work}
\label{sec:deadends}

This section is longer than most ``negative results'' sections because
several of the failures were more informative than the successes. The
architecture we ended up with --- RCB, Shamir's trick, bytecode
interpreter, $q=t[\text{top}]$ --- was not a foregone conclusion. We
built prototypes of three alternative architectures to find out whether
it was the right one. It was.

\subsection{Implementation-level failures}

\begin{description}[leftmargin=1em,style=nextline]
\item[Merging schoolbook and reduce into one loop.] The multiply
  loop and the reduce loop have similar shapes --- walk an array,
  multiply by something, accumulate. Merging them would share the
  loop machinery. It also overflows: the accumulated carry during
  the multiply phase can exceed what the reduce phase expects, and
  you get wrong answers on about one input in $2^{30}$. Reverted
  after a Wycheproof failure.
\item[Zero-skip in the scalar loop.] Half the scalar bits are zero,
  so skipping the add on zero bits should be a speedup. It was 40\%
  \emph{slower}: the branch mispredicts roughly half the time (the
  bits are effectively random) and the penalty swamps the saved
  work. \texttt{perf stat} showing branch-miss rate was the
  diagnostic that killed this idea.
\item[Solinas reduction for both moduli.] The prime $p$ has a
  beautiful sparse form ($2^{256} - 2^{224} + 2^{192} + 2^{96} - 1$)
  that admits a reduction with no multiplies at all, just shifts and
  adds. But $n$ doesn't have that form, so you need \emph{two}
  reducers instead of one. For a speed build that's worth it
  (\texttt{-DSOLINAS\_P} exists); for the size floor, one generic
  reducer serving both moduli is smaller.
\end{description}

\subsection{Dropping Shamir's trick}

The scalar multiplication loop computes $u_1 \cdot G + u_2 \cdot Q$.
Shamir's trick (\S\ref{sec:slotshift}) does this in one pass: at each
bit, double once, then conditionally add $G$ or $Q$ or both. The
alternative is simpler: compute $u_1 \cdot G$, compute $u_2 \cdot Q$,
add the results. One routine called twice instead of a two-scalar
loop with its backup tables and dual bit-tests.

This seemed promising for size. The reference implementation
\texttt{p256-m}~\cite{p256m} does exactly this, and its README says
``in order to minimize code size.'' There is also a neat theorem
(Hamburg~\cite{hamburg}) showing that if you prepare the scalar
carefully --- force it odd, initialise the ladder to a specific state
--- then inside the loop the running point can never coincide with the
base point or its negative. That means the edge cases RCB exists to
handle (\S\ref{sec:rcb}) never occur, and you can use an older,
simpler, faster addition formula with no branches at all.

We built it. It measured \spent{277}. Here is where the bytes went:

\begin{center}
\begin{tabular}{@{}lrl@{}}
\toprule
Component & $\Delta$ & Why it costs \\
\midrule
Bytecode for the two formulas & \saved{7} & About what you'd expect \\
Conditional ladder setup       & \spent{50} & Bytecode can't branch \\
Two-call verify tail           & \spent{117} & Arg setup $\times$ 2; combining add \\
Coordinate-system mismatch     & \spent{33} & Needs a mod-$p$ inversion back \\
\bottomrule
\end{tabular}
\end{center}

The bytecode saving is real --- RCB's 43-operation complete formula
compresses to \bytes{88}, and the two simpler incomplete formulas
together compress to \bytes{80}. Eight bytes is the cost of
completeness at the bytecode level. That is astonishingly cheap.

Everything else is infrastructure we already have for free and would
have to rebuild. The ladder setup needs a branch (is the scalar odd or
even?), but the bytecode interpreter is straight-line --- so the
branch lives in assembly, and each arm is a full setup sequence. The
two results need combining with a final add, but that add can
hit the edge cases the ladder avoided --- so the three-way branch
comes back for one call. And the simpler formulas use a different
coordinate system (Jacobian) than ours (homogeneous), so converting
for the final combine brings back the mod-$p$ inversion we deleted in
\S\ref{sec:projcheck}.

The lesson is not that \texttt{p256-m}'s authors were wrong. Their
baseline is C with separate functions and a branchy verify tail; adding
a three-way combine costs them twenty bytes. Our baseline is 27 bytes
of fall-through. The same structural change has wildly different cost
against different baselines.

\subsection{Single-kernel modular multiply}

The modular multiply (\S\ref{sec:nomont}) interleaves product and
reduction row by row. Gueron~\cite{gueron} notes, for the general
Montgomery case, that you can deinterleave: compute the full product,
then the full reduction, using one big multiply routine called twice.
The appeal for size is code reuse --- one kernel, two call sites.

For P-256 specifically there seemed to be a bonus. Montgomery
reduction computes a correction factor by multiplying by a magic
constant, $-p^{-1}$ mod $2^W$ where $W$ is the limb width. For $W \le
96$, $p$'s low bits are all ones, so this constant is exactly 1 and
the multiply is free. If we take $W = 256$ --- one giant limb --- would
the constant still be 1, making the correction step a no-op?

It would not:
\[
-p^{-1} \bmod 2^{256} = \texttt{0xffffffff\,00000002}\cdots\texttt{00000001}
\]
The 96-bit all-ones run stops at bit 96. Beyond that, $p$ has
structure, and its inverse has structure too. The free correction
disappears.

The deeper problem: our inner loop already is the factored
form. One small routine (\texttt{.Lcnt}, 15 bytes) is called twice per
row --- once for product, once for reduction. Deinterleaving would mean
two outer loops instead of one, duplicating the loop counter and tail
machinery. We know this costs bytes because an earlier commit went the
other direction, merging two separate loops into the current
interleaved form, and measured \saved{7}.

\subsection{Scalar splitting and sparse recodings}

Two standard speedups for scalar multiplication:
\begin{itemize}
\item Some curves have a cheap ``endomorphism'' --- a way to multiply
  any point by a fixed large constant almost for free --- and you can
  split a 256-bit scalar into two 128-bit halves, halving the loop.
  P-256 has no such endomorphism. This is a property of the curve's
  algebraic structure (a number called the $j$-invariant that
  classifies curves by their endomorphism ring) and is not fixable.
\item Recoding the scalar in a sparser base (non-adjacent form, window
  methods) reduces the number of additions, because more digits are
  zero. This saves arithmetic operations --- but our bytecode is two
  bytes per operation whether you execute it once or a thousand times.
  Fewer \emph{executions} of the add step saves time, not code. The
  recoder itself is pure overhead.
\end{itemize}
Both are genuine speedups; neither touches the size axis in the
direction we want.

\subsection{Finding bytecode hidden in other bytes}

The bytecode is 2-byte operations, and we have hundreds of bytes of
constants and instruction encodings sitting in the binary. Could any
of those bytes, by accident, already encode a useful bytecode
sequence? If three consecutive operations from our add formula happened
to appear inside the bytes of a curve constant, we could point the
interpreter there and save six bytes of storage.

We scanned exhaustively. Nothing matches. The bytecode encoding is
dense --- roughly two-thirds of all byte values are valid first bytes
of \emph{some} operation --- which means random bytes produce long runs
of syntactically valid but semantically useless operations. There is a
16-operation ``valid'' run inside one of the constants; it does
nothing useful.

The same question one level down: x86 instructions are
variable-length, so the bytes of one instruction, starting at an
offset, decode as a different instruction. If a useful instruction is
hidden inside another, you can jump to the offset and get it for free.
We found exactly one such gadget: inside the four-byte encoding of
\texttt{imul eax, r11d}, starting at byte three, is \texttt{scasd;
ret} --- a legitimate two-instruction function. It advances \texttt{rdi}
by four. We work in eight-byte words. Useless.

A tamer version did find something: two functions that end with the
same three bytes (\texttt{stosq; ret}) and sit 43 bytes apart. One
jumps to the other's tail. \saved{1}.

% ======================================================================
\section{Other limb widths: 11$\times$24, 5$\times$54, 5$\times$56}
\label{sec:limbs}

Everything above is one architecture: eight 32-bit words per field
element, plain (non-Montgomery) reduction. That choice was made early
and then optimised hard. But it's a choice, and other choices carve
up the problem differently. This section sketches the alternatives
explored in parallel --- none has beaten \bytes{890}, but each
taught something about where the bytes actually go.

\subsection{What a limb is and why its width matters}

A 256-bit integer doesn't fit in a general-purpose register --- the
64-bit kind you can \texttt{add} and \texttt{mul} --- so it's stored
as an array of $K$ smaller pieces --- \emph{limbs} --- of $W$ bits
each, with $K \cdot W \ge 256$. The slack $K \cdot W - 256$ is
headroom for carries to accumulate between reductions --- a style
called \emph{lazy carrying}. The main build uses $K=8$, $W=32$: no
slack, carries propagated immediately.

Multiplying two $K$-limb numbers by the schoolbook method is $K^2$
single-limb multiplies --- for $K=8$ that's 64; for $K=5$ it's 25.
Fewer limbs means dramatically less arithmetic. But fewer limbs means
wider limbs, and wider limbs have wider products: a $54 \times 54$
bit multiply overflows a 64-bit register and needs 128-bit
intermediate results, which on x86-64 means the two-register
\texttt{rdx:rax} output of \texttt{mul} and extra bookkeeping to
handle it. A $24 \times 24$ bit multiply fits comfortably in 64 bits
with room to accumulate a column's worth of partial products before
anything overflows.

So the trade-off runs: small $K$ saves multiplications but pays for
wide-product handling; large $K$ keeps each step simple but does
many more of them. Where the minimum lands depends on the
instruction set, the reduction strategy, and --- it turns out --- a
bit-level accident in the constant $p$.

\subsection{Quotient estimation, in plain terms}

Before getting to the accident, it helps to see what all reduction
strategies are really doing. After multiplying two numbers modulo
$p$, you have a result up to $p^2$ --- roughly twice as many bits as
you want --- and need to bring it back below $p$. Conceptually this
is division: compute $t \bmod p$ by finding how many times $p$ goes
into $t$ and subtracting that many copies.

Division is what you want to avoid. It's the slowest arithmetic
instruction on almost every chip ever made, and it doesn't pipeline.
Every practical reduction scheme replaces the division with a
\emph{guess}: estimate the quotient cheaply, subtract, see how close
you got, correct if needed. The schemes differ in how they guess.

The $q=t[\text{top}]$ reduce (\S\ref{sec:nomont}) guesses by reading
off the top word --- which works because $p$'s top word is all ones,
so $p$ is within a hair of the next power of two and ``how many
times does $p$ go in'' is almost exactly ``how many times does
$2^{256}$ go in,'' which is just a shift. It's the same reason long
division by hand is easy when the divisor is 999: you barely have to
think about the trial digit.

Montgomery reduction~\cite{montgomery} guesses from the other end.
It asks: what multiple of $p$ can I add to $t$ so that the
\emph{bottom} $W$ bits become zero? Then those zero bits are shifted
off, and the value shrinks by $W$ bits per round. Finding that
multiple normally takes one extra multiply per round --- you compute
$q = (t \bmod 2^W) \cdot c$, where the constant $c = -p^{-1} \bmod
2^W$ is precomputed once and baked into the binary. One multiply by a
fixed constant per limb, which is why the main build dropped
Montgomery: that multiply and its constant are overhead.

\subsection{Ninety-six ones in a row}

Here is the accident. Recall $p = 2^{256} - 2^{224} + 2^{192} +
2^{96} - 1$. That trailing $-1$ borrows all the way up to bit 96
before the $+2^{96}$ term stops it --- so bits 0 through 95 of $p$
are all ones. You can read it straight off the hex from
\S\ref{sec:nomont}: the bottom three words are
\texttt{ffffffff\,ffffffff\,ffffffff}.

What does this do to the Montgomery constant $c = -p^{-1} \bmod 2^W$?
If $W \le 96$, then $p \bmod 2^W$ is a $W$-bit block of ones --- that
is, $p \equiv -1 \pmod{2^W}$. Then $p^{-1} \equiv -1$ as well (since
$-1$ is its own inverse), and $c = -(-1) = 1$.

\textbf{The per-limb Montgomery multiply is by 1.} It disappears. For
any limb width up to 96 bits, Montgomery reduction mod $p$ costs no
more per round than the $q=t[\text{top}]$ scheme does --- the guess
is free in both. Every limb width on the table except $8 \times 32$
in its non-Montgomery form exploits this: the Montgomery machinery
that \S\ref{sec:nomont} threw out comes back, but without the part
that made it expensive.

(The order $n$ doesn't share this structure --- its low bits are
\texttt{fc632551} --- so the mod-$n$ inverter still pays for its
Montgomery constant. But inversion runs once; the field multiplier
runs thousands of times.)

\subsection{Proving the limbs don't overflow}

A limb representation can be correct on every test vector and still
be wrong. The failure mode: some intermediate value overflows its
accumulator on roughly one input in $2^{30}$, and your test suite
never hits it. We know this happens because it happened
(\S\ref{sec:deadends}).

So for each $(K, W)$ pair we prove the bounds symbolically. The RCB
formula (\S\ref{sec:rcb}) is a fixed sequence of 43 multiplies and
adds; with lazy carrying, each add can grow its output's limbs
beyond $W$ bits, and the question is whether any limb ever exceeds
what the multiplier can handle. For $K = 11$, $W = 24$, a 64-bit
accumulator holds $K$ products of $29 \times 29$ bits with room to
spare --- but only if the inputs really are 29 bits. Tracing the
formula operation by operation, the worst case is a $\times 12$
growth chain that lands the output at about $2^{258}$, or four times
$p$: the top limb reaches 18 bits, the rest stay under 29. At $K
\cdot W = 264$ bits of capacity, this converges; at 260 it would not.

The proof is a few dozen lines of Python that walks the formula and
tracks interval bounds. It runs in the build. A configuration that
doesn't converge doesn't compile.

\subsection{The tracks}

\begin{center}
\begin{tabular}{@{}lccrl@{}}
\toprule
Track & $K$ & $W$ & Size & Reduction strategy \\
\midrule
$8 \times 32$  & 8  & 32 & \bytes{890}  & non-Montgomery, $q=t[\text{top}]$ \\
$11 \times 24$ & 11 & 24 & \bytes{1068} & Montgomery, $c_p = 1$, 64-bit accumulators \\
$5 \times 56$  & 5  & 56 & \bytes{1084} & Montgomery, $c_p = 1$, byte-aligned decode \\
$5 \times 54$  & 5  & 54 & \bytes{1097} & Montgomery, $c_p = 1$, 128-bit products \\
\bottomrule
\end{tabular}
\end{center}

\paragraph{$11 \times 24$.} Eleven limbs of 24 bits gives 264 bits of
capacity --- 8 bits of slack per field element. A $24 \times 24$
product is 48 bits; even after summing eleven of them down a column,
the total stays under 64 bits. No \texttt{rdx:rax} splitting, no
128-bit adds, just plain \texttt{imul} and \texttt{add}. The
multiplier inner loop is the simplest of any track. The cost:
$11^2 = 121$ multiplies per field operation versus $8^2 = 64$, almost
double, and the carry-propagation code that runs after each operation
is longer because there are more limbs to walk.

\paragraph{$5 \times 54$.} Twenty-five multiplies. The $K^2$ scaling
wins decisively on arithmetic count. But 54-bit limbs mean 108-bit
products, handled as \texttt{rdx:rax} pairs, and the reduction
involves 128-bit-wide shift-and-add chains. And 54 is not a multiple of 8:
decoding a big-endian input into 54-bit limbs requires loading a
word, shifting by a variable amount, and masking --- a bit-twiddling
loop that the byte-granular tracks don't need.

\paragraph{$5 \times 56$.} Same $K=5$ as above, so the same 25
multiplies and the same $c_p = 1$. The difference is that 56 bits is
exactly 7 bytes. Decoding a 56-bit limb from a byte buffer is an
unaligned 8-byte load followed by masking off the top byte --- no variable shift
loop, no bit-offset bookkeeping. Two extra bits of width per limb
buys back a chunk of the parsing machinery for nearly nothing, and
$56 < 96$ so the Montgomery constant is still 1.

It does save: $5 \times 56$ finished \bytes{13} smaller than
$5 \times 54$. The byte-aligned decode is worth more than the
slightly wider products cost. The $5 \times 54$ track has one
consolation prize: because its bytecode pointer happens to live in
\texttt{rsi} rather than \texttt{rbx}, the \texttt{rbx} register is
free to hold the jump-table base, enabling a one-byte instruction
called \texttt{xlatb} (a leftover from 8-bit days: it does
\texttt{al} $\leftarrow$ \texttt{[rbx+al]}, a table lookup in one
byte). That
trick ports to $5 \times 56$ too, but only after moving the bytecode
pointer --- and $5 \times 56$ already had it moved, for a different
reason, so the net was larger there. Tricks compound with
architectural accidents.

\subsection{Why $8 \times 32$ is still ahead}

The $8 \times 32$ track has one structural advantage none of the
others share: a single reducer serves both moduli. Both $p$ and $n$
have a top 32-bit word of \texttt{0xFFFFFFFF}, so the
$q=t[\text{top}]$ scheme handles both with one code path. Every
Montgomery track needs the $c_p = 1$ trick for the field and
something else for the order --- two reducers, or one reducer with a
mode switch.

All four tracks now have roughly the same tricks applied --- the
encoding wins cross-pollinated freely once discovered --- and the gap
held at about \bytes{180}. The two-reducer cost is real. The
Montgomery tracks did find things the $8 \times 32$ build cannot use:
cancelling $R$-factors in the projective check (\saved{27}), a
signed-limb representation of $p$ built from six increment/decrement
instructions (\saved{4}). Not enough to close the gap, but enough to
make the comparison fair.

The Montgomery tracks are faster. At the same size, they would win on
both axes. They are not the same size.

\subsection{The non-portability of tricks}

The $11 \times 24$ track found a neat consolidation: merge two loops
(``add the modulus while the value is negative'' and ``subtract the
modulus until it goes negative'') into one, using the CPU's
sign-extend instruction to pick the direction each iteration. The
termination test is a single \texttt{or edx, eax; js} --- if we just
subtracted, the \texttt{or} is negative regardless (loop); if we just
added, the \texttt{or} equals the top limb, and the sign flag tracks
it (exit when non-negative). \saved{3}.

Porting to the other two Montgomery tracks saved only one byte, or
nothing. Two costs materialised that $11 \times 24$ didn't have.
First, its carry-propagation routine happened to end with the top limb
already in a register --- the trick got its input for free. The other
tracks' routines don't end that way and need an extra load. Second,
eleven 24-bit limbs means the top limb after carry propagation fits in
18 bits, so bit 31 of a 32-bit register is genuinely the sign bit and
a 2-byte \texttt{or} works. The five-limb tracks have wider top limbs
--- up to 44 bits --- and bit 31 is data, not sign. They need the
3-byte wider-register \texttt{or}.

The trick was not ``merge the loops.'' It was ``merge the loops when
your carry-prop tail happens to leave the operand loaded and your limb
width happens to fit in half a register.'' Architectural coincidences
compound; so does their absence.

% ======================================================================
\section{General principles}

If you're trying to shrink your own binary, here's what transfers:

\begin{enumerate}
\item \textbf{Count call sites, not functions.} A function called
  once is as expensive as its call site plus its body. A function
  called sixty times is sixty call sites plus one body. The
  interpreter trick (\S\ref{sec:bytecode}) wins because it collapses
  sixty 15-byte call sites into sixty 2-byte data entries.

\item \textbf{Algebra beats branching.} The projective check
  (\S\ref{sec:projcheck}) and the RCB formula (\S\ref{sec:rcb}) both
  replace control flow with arithmetic. Arithmetic is data;
  branches are code. Data is cheaper.

\item \textbf{Your constants are not random.} The \texttt{bt}-on-$n$
  trick (\S\ref{sec:btcn}), the $q=t[\text{top}]$ reduce
  (\S\ref{sec:nomont}), and building $p$ at runtime
  (\S\ref{sec:runtimecp}) all exploit specific bit patterns in the
  P-256 constants. Stare at your constants in hex. Stare at them in
  binary. There is structure there.

\item \textbf{Layout is a degree of freedom.} Slot numbering
  (\S\ref{sec:slotshift}), function order (short jumps), constant
  order (block copies), rodata-before-text (push imm8) --- none of
  these affect semantics, all of them affect size. Change them last,
  after the architecture settles, because they're brittle.

\item \textbf{Track both axes.} It's easy to micro-optimise yourself
  into a corner where you've saved 5 bytes and lost $10\times$ speed.
  Plot size against cycles for every checkpoint; if your new point is
  dominated by an old one on both axes, you've wasted your time.

\item \textbf{Tests before bench.} Every single one of the $\sim$130
  commits in this project passed 607/607 test vectors before being
  committed. Several ideas were killed by tests (see above). A
  smaller verifier that accepts forgeries is negative progress.
\end{enumerate}

% ======================================================================
\begin{thebibliography}{9}

\bibitem{fips186}
National Institute of Standards and Technology.
\textit{Digital Signature Standard (DSS)}.
FIPS PUB 186-5, February 2023.
\url{https://doi.org/10.6028/NIST.FIPS.186-5}

\bibitem{rcb}
Joost Renes, Craig Costello, and Lejla Batina.
\textit{Complete addition formulas for prime order elliptic curves}.
EUROCRYPT 2016. IACR ePrint 2015/1060.
\url{https://eprint.iacr.org/2015/1060}

\bibitem{efd}
Daniel J.\ Bernstein and Tanja Lange.
\textit{Explicit-Formulas Database}.
\url{https://hyperelliptic.org/EFD/}

\bibitem{montgomery}
Peter L.\ Montgomery.
\textit{Modular multiplication without trial division}.
Mathematics of Computation, 44(170):519--521, 1985.

\bibitem{wycheproof}
Google Security Team (now community-maintained under C2SP).
\textit{Project Wycheproof: Testing crypto libraries against known attacks}.
\url{https://github.com/C2SP/wycheproof}

\bibitem{hac}
Alfred J.\ Menezes, Paul C.\ van Oorschot, and Scott A.\ Vanstone.
\textit{Handbook of Applied Cryptography}.
CRC Press, 1996.
\url{https://cacr.uwaterloo.ca/hac/}

\bibitem{straus}
Ernst G.\ Straus.
\textit{Addition chains of vectors (Problem 5125)}.
American Mathematical Monthly, 71(7):806--808, 1964.

\bibitem{agner}
Agner Fog.
\textit{Instruction tables: Lists of instruction latencies,
throughputs and micro-operation breakdowns}.
\url{https://www.agner.org/optimize/instruction_tables.pdf}

\bibitem{hamburg}
Mike Hamburg.
\textit{Faster Montgomery and double-add ladders for short Weierstrass curves}.
IACR Transactions on Cryptographic Hardware and Embedded Systems
2020(4):189--208. IACR ePrint 2020/437.
\url{https://eprint.iacr.org/2020/437}

\bibitem{p256m}
Manuel P\'egouri\'e-Gonnard.
\textit{p256-m: Minimal P-256 ECDSA / ECDH in portable C}.
\url{https://github.com/mpg/p256-m}

\bibitem{gueron}
Shay Gueron.
\textit{Efficient software implementations of modular exponentiation}.
IACR ePrint 2011/239.
\url{https://eprint.iacr.org/2011/239}

\end{thebibliography}

\end{document}
